\section{Rotation Relaxation and a Time Bound}
\label{sec:bounds}

% 证据：prior/.../超长周期_数学/第三轮_定理与证据.md（命题 1、2(i)，定理 8，引理 7）；研究笔记/06（接入定理、小例子）；
%       code/t6_job.py（速率 LP，相位层面）、prior/.../多型号多地/freshlb.py（时间展开 LP）；code/theory/toy_connector.py（例子穷举）。
% 不能写：训练实例上的数字；"最优解一定是接入形式"（未证）；多机型的"最多差 1 台"（只在单机型上有机证）。

This section answers two questions.
How many resources does MDRR need at least?
And how long a start-up must a schedule optimize before it can simply repeat a rotation?
Both answers come from the maintenance cycle of a robot, which we call its \emph{rotation}.

\subsection{Rate Relaxation}
\label{sec:bounds:rate}

Let $\tau_{k,i,s}$ be the number of slots between two ready times of a type-$k$ robot that serves site $s$ in phase $i$; by \eqref{eq:step}, $\tau_{k,i,s} = \lceil \delta_s/u \rceil + \lceil (o_{k,i}+\delta_s)/u \rceil$.
Let $x_{k,i,s} \ge 0$ be the long-run number of phase-$i$ duties per slot that type-$k$ robots serve at site $s$.
The \emph{rate relaxation} is
\begin{align}
  R = \min \;& \textstyle\sum_k N_k + \sum_j M_j \label{eq:rate}\\
  \text{s.t. }& \textstyle\sum_{k,i} d_k\, x_{k,i,s} = K_s, && s \in \mathcal{S}, \nonumber\\
  & \textstyle\sum_s x_{k,i,s} = \sum_s x_{k,i+1,s}, && k \in \mathcal{K},\ i < L_k - 1, \nonumber\\
  & \textstyle\sum_{i,s} \tau_{k,i,s}\, x_{k,i,s} \le N_k, && k \in \mathcal{K}, \nonumber
\end{align}
together with the analogous tool constraints: each duty assigns $q_k$ tool units, and a tool of type $j$ is busy for $\alpha_{k,s} + d_k + \beta_{k,s,j}$ slots per duty.

\begin{proposition}
\label{prop:rate}
Every feasible schedule uses at least $\lceil R \rceil$ resources.
\end{proposition}

\begin{proof}
Count over a window $[t, t + T)$ after the start-up window.
Site $s$ needs $K_s T$ duty slots.
A type-$k$ robot that serves $n_{i,s}$ phase-$i$ duties at site $s$ in the window is busy for at least $\sum_{i,s} \tau_{k,i,s} n_{i,s} - O(1)$ slots, and it serves the phases in cyclic order, so its phase counts differ by at most one.
Divide all counts by $T$ and let $T \to \infty$.
The limits satisfy the constraints of \eqref{eq:rate}, so the fleet is at least $R$.
Tools are handled in the same way.
\end{proof}

The rate relaxation ignores how the fleet starts.
It is a lower bound for any schedule, but it is reached only if the robots are already spread over their cycles.

\subsection{Start-Up Relaxation}
\label{sec:bounds:startup}

A new fleet is not spread over its cycles.
All robots are in phase $0$ at time $0$, so their first overhauls come close together.
To see this effect we expand time over a horizon $[0, H)$.
Each robot is a path in a graph whose nodes are (type, phase, ready slot) and whose arcs are waits, duties at eligible sites by \eqref{eq:step}, and early overhauls inside $[0, W)$.
Every site needs $K_s$ robots on duty in each open slot of $[0, H)$; arcs that end after $H$ are free.
Tools are paths in a similar graph and are coupled to the duties they serve.
Let $R_0(H)$ be the value of the linear program over these path flows.

\begin{proposition}
\label{prop:startup}
For every $H$, every feasible schedule uses at least $\lceil R_0(H) \rceil$ resources, and $R_0(H)$ is nondecreasing in $H$.
\end{proposition}

\begin{proof}
Cut a feasible schedule at $H$.
Its robots and tools give integral flows that satisfy all constraints on $[0, H)$.
A longer horizon only adds constraints.
\end{proof}

The dual of this linear program puts a price on every (site, slot).
A price is high where many robots are away for their first overhaul at the same time.
These prices show where in time the start-up is expensive, and we use them in Section~\ref{sec:experiments}.

\subsection{One Robot Type: What the Bounds Mean}
\label{sec:bounds:single}

The following two results concern one robot type and $K$ posts at one site, with time counted in duties.
After a duty in phase $i$, a robot must wait $g_i$ duties before its next one.
Its cycle has length $S = \sum_i g_i$, and the rate bound is $\lceil K S / L \rceil$.

\begin{proposition}[Even rotation reaches the rate bound]
\label{prop:orbit}
If $K = mL$ for an integer $m$, then $mS$ robots suffice, which equals the rate bound.
\end{proposition}

\begin{proof}
Let $A = \{0, g_0, g_0 + g_1, \dots\}$ be the duty times of one cycle, taken modulo $S$.
Put $m$ robots at every offset $b \in \{0, \dots, S-1\}$ and let each repeat its cycle from $b$ without waiting.
Slot $x$ is served by the robots with $x - b \in A \pmod S$.
There are exactly $L$ such offsets, so $mL = K$ robots are on duty in every slot.
\end{proof}

Without early overhaul, a new fleet can need many more robots than the rotation.

\begin{proposition}[Start-up cost can grow with the number of posts]
\label{prop:linear}
Let $g = (1, 2k)$, so $S = 2k + 1$.
Without early overhaul, a new fleet needs at least $K (S^2 + 1)/(2S)$ robots, while $\lceil K S / 2 \rceil$ robots suffice if the robots may start anywhere in their cycles.
The difference grows linearly in $K$.
\end{proposition}

\begin{proof}[Proof sketch]
Give slot $t$ the weight $w(t) = \lfloor t/2 \rfloor + 1$ for $t < 2k$, $w(2k) = k + 1$, $w(t) = \lceil (4k - t)/2 \rceil$ for $2k < t < 4k$, and $0$ otherwise.
A new robot serves its first four duties at $t_1 < t_2 < t_3 < t_4$ with $t_2 \ge t_1 + 1$, $t_3 \ge t_2 + 2k$ and $t_4 \ge t_3 + 1$.
A case check shows that the weights it collects sum to at most $S$.
Every slot has $K$ robots on duty, so all robots together collect $K \sum_t w(t) = K(2k^2 + 2k + 1)$, and at least $K (2k^2+2k+1)/S = K(S^2+1)/(2S)$ robots are needed.
For the second claim, even $K$ is Proposition~\ref{prop:orbit}.
For odd $K$, serve $K-1$ posts as in Proposition~\ref{prop:orbit} and the last post with $(S+1)/2$ robots that wait one extra slot after each overhaul: their cycle $(1, 2k+1)$ has even length $S + 1$, and robots at offsets $0, 2, \dots, S-1$ cover every slot once.
\end{proof}

The weights in this proof are a dual solution of the start-up relaxation.
They are large in the first $2S$ slots and zero afterwards: the extra cost of a new fleet is paid at the start.

\subsection{A Time Bound: Entering a Rotation}
\label{sec:time}

The natural period of a good schedule can be very long, because the cycles of different types do not divide each other.
Encoding one period is therefore not an option.
The next result shows that this is not needed: any rotation can be entered from a fresh start within one maintenance interval.

Let $\sigma$ be a feasible periodic schedule in which robots may start anywhere in their cycles.
It uses $N_\sigma$ robots and $M_\sigma$ tools.
Let $C_\sigma$ be the longest time, over all robots of $\sigma$, between two consecutive returns to phase $0$.

\begin{theorem}[Entering a rotation]
\label{thm:enter}
Let $T_0 = c_0 + C_\sigma$.
There is a fresh-start schedule that agrees with $\sigma$ in every slot $t \ge T_0$ and uses $N_\sigma + E$ robots and $M_\sigma + F$ tools.
Here $E$ and $F$ are the resources needed to serve, from a fresh start, the duties of $\sigma$ in $[0, T_0)$ that the copied robots below do not serve.
\end{theorem}

\begin{proof}
For a robot $r$ of $\sigma$, let $b_r$ be its last return to phase $0$ before $T_0$; then $b_r \ge T_0 - C_\sigma \ge c_0$.
Give $r$ a new robot $r'$.
A new robot is in the same state as a robot that has just finished an overhaul, so $r'$ can wait in phase $0$ until $b_r$ and then repeat every action of $r$.
Give every tool of $\sigma$ a new tool that repeats its duties when the robot of that duty has started copying, and stays idle otherwise.
From $T_0$ on, every $r'$ copies $r$, so the schedule equals $\sigma$.
Before $T_0$, the copies serve all duties of $\sigma$ except those of robots that have not started copying.
The $E$ robots and $F$ tools serve these duties.
No two copies conflict, because the robots and tools of $\sigma$ do not.
\end{proof}

The theorem bounds the start-up in time, and the bound does not depend on the number of sites, posts or robot types.
It does not bound $E$.
Without early overhaul, $E$ can grow with the number of posts, as Proposition~\ref{prop:linear} shows.
Early overhaul lets the copies serve the missing duties themselves: before $b_r$, robot $r'$ may serve a few duties and then go to overhaul, as long as it is back in phase $0$ by $b_r$.
Fig.~\ref{fig:enter} shows a case where this removes all extra robots.
For one robot type the following condition makes the extra robots unnecessary.

\begin{proposition}[Self-cover]
\label{prop:selfcover}
Consider one robot type and $K = mL$ posts at one site, with time counted in duties and early overhaul allowed in the start-up.
If $g_p \le g_{i+p}$ for all $1 \le i \le L-1$ and $0 \le p \le L-2-i$, then a new fleet of $KS/L$ robots, the rate bound, enters the even rotation of Proposition~\ref{prop:orbit} within $S$ slots.
In particular this holds if the gaps before the overhaul are nondecreasing, and always if $L = 2$.
\end{proposition}

\begin{proof}
Let $s_i = g_0 + \dots + g_{i-1}$.
In the even rotation, a robot with offset $b \in \{0, \dots, S-1\}$ is in phase $0$ at slot $b$; its copy enters there.
Before $b$, the rotation still has this robot's previous cycle: its phase-$i$ duties at slots $t_i = b - S + s_i$ for $i \ge i^*$, where $i^* \ge 1$ is the smallest index with $s_i \ge S - b$.
Let the new robot serve these duties itself, as its phases $0, 1, \dots$.
Consecutive duties are $g_{i^*+p}$ slots apart, and the robot needs only $g_p \le g_{i^*+p}$ after its phase $p$.
Its last such duty is at $b - g_{L-1}$, in phase $L - 1 - i^* \ne L-1$, so it may go to overhaul and is back in phase $0$ at $b$.
Every slot is then served exactly as in the rotation.
\end{proof}

The example of Fig.~\ref{fig:enter} is the case $L = 2$.
An exhaustive search confirms there that without early overhaul a new fleet needs $7$ robots instead of $6$.

\begin{figure}[t]
  \centering
  \figplaceholder{3.6cm}{Fig.~X (to draw): $g = (2, 4)$, two posts. Top: the even rotation with 6 robots (offsets 0--5). Bottom: a fresh fleet copies it; robots 4 and 5 serve slots 0 and 1, go to early overhaul, and join the rotation at slots 4 and 5. Without early overhaul 7 robots are needed (exhaustive check).}
  \caption{Entering a rotation from a fresh start. One cycle of 6 slots, two duties per cycle, two posts. Early overhaul lets the fleet of the rotation fill the start-up gap.}
  \label{fig:enter}
\end{figure}

\begin{corollary}[A short window suffices to represent the schedule]
\label{cor:window}
For every rotation $\sigma$, the schedule of Theorem~\ref{thm:enter} differs from $\sigma$ only in $[0, c_0 + C_\sigma)$.
An encoding that fixes the schedule after this window to a rotation, and searches only the decisions inside it, can therefore represent it.
\end{corollary}

Corollary~\ref{cor:window} is the reason our method encodes only a short start-up window and replays a rule afterwards.
How short the window can be in practice is measured in Section~\ref{sec:experiments}.
